\documentclass{configs/idp-model}
\usepackage{float}
\usepackage{tikz}
\usetikzlibrary{positioning}
\usepackage{graphicx}
\usepackage{caption}
\usepackage{longtable}
\usepackage{array}


%PREENCHER TODOS OS CAMPOS ABAIXO
\autor{Gabriel Queiroz Teles}


% Utilizar um dos dois cursos abaixo
\cienciadacomputacao
%\engenhariadesoftware


\titulo{Autonomous Object Capture System by UAVs Using Visual Servoing}
\ano{2026}

% Usar orientador ou orientadora
% - Argumento opcional caso o orientador seja de uma instituição diferente
\orientador[IDP]{Patrícia da Silva Oliveira}
%\orientadora{Dra. Ada Lovelace}

% Para membros da banca, são três argumentos obrigatórios:
% 1 - Nome do membro avaliador
% 2 - Se o mesmo é membro interno ou externo à instituição
% 3 - Instituição a qual o membro está vinculado
\membrobancai{Dr. Dom Quixote de la Mancha}{Internal Member}{IDP}
\membrobancaii{Dra. Anna Karienina}{Internal Member}{IDP}

\palavraschave{ Veículos Aéreos Não Tripulados. Servovisão. YOLO. Filtro de Kalman. Manipulação Aérea.}
\keywords{Unmanned Aerial Vehicles. Visual Servoing. YOLO. Kalman Filter. Aerial Manipulation.}

\dataaprovacao{08/05/2025}



\begin{document}

\setcounter{secnumdepth}{3} 
\setcounter{tocdepth}{3}

% Não alterar
\pagenumbering{roman}
\capaprincipal
\folharosto
\folhaaprovacao

% Opcional
\dedicatoria
\agradecimentos


% Não alterar 
% Resumo e abstract - alterar os arquivos na pasta "partes" NÃO ALTERAR O NOME DELES
\abstract
\resumo

%List of Figures e Tables são optativos. Caso não esteja utilizando no seu projeto, comentar as linhas abaixo
\listoffigures
\listoftables

\tableofcontents
\pagenumbering{arabic}


% Para as seções, o comando \secao deve ter a forma \secao{título}{arquivo .tex}

% Os comandos \subsecao e \subsubsecao seguem a forma \comando{Título}, sem a necessidade de passar um arquivo .tex

\secao{Introduction}{capitulos/introduction}
\secao{Literature Review}{capitulos/literature-review}
\secao{Methodology}{capitulos/methodology}
\secao{Results and Discussion}{capitulos/results}
\secao{Conclusion}{capitulos/conclusion}

% Não alterar
\referencias

% Opcionais
\apendice[ap:meuapendice]{Título do apêndice A}{capitulos/apendicei}
% Alternativa sem o primeiro parâmetro também funciona
% \apendice{Título do apêndice A}{partes/apendicei}

\anexo{Título do Anexo A}{capitulos/anexoi}


% Não alterar
\capafinal

 
\end{document}